\documentclass[a4paper,12pt] {article}

\usepackage[utf8]{inputenc}
\usepackage[english,russian]{babel}

\usepackage{geometry}
\geometry{left=3cm}
\geometry{right=1cm}
\geometry{top=2cm}
\geometry{bottom=2cm}
\linespread{1.3}

\def\htag#1{\#\textit{#1}}  % short alias for hashtags style

\usepackage{graphicx}
\graphicspath{{pic/}}
\DeclareGraphicsExtensions{.png,.jpg}

\newcommand{\tocsecindent}{\hspace{7mm}}


\begin{document}


\begin{titlepage}
		
\thispagestyle{empty}
		
  \begin{center}
  
   
   ПРАВИТЕЛЬСТВО РОССИЙСКОЙ ФЕДЕРАЦИИ \\
   
   \vspace{0.25cm}
   ФЕДЕРАЛЬНОЕ ГОСУДАРСТВЕННОЕ БЮДЖЕТНОЕ
   ОБРАЗОВАТЕЛЬНОЕ УЧРЕЖДЕНИЕ ВЫСШЕГО ОБРАЗОВАНИЯ 
   «САНКТ-ПЕТЕРБУРГСКИЙ ГОСУДАРСТВЕННЫЙ УНИВЕРСИТЕТ» 
   (СПбГУ)
   
   \vspace{0.25cm}
   Направление: Прикладные математика и физика
   
 
   \begin{figure}[h!]
   	\begin{center}
   		\includegraphics[width=0.25\linewidth]{Gerb}
   	\end{center}
   \end{figure}


     \textbf{\Large Полуавтоматическое структурирование
     изображений в социальной сети с помощью 
     методов машинного обучения}
     
 \bigskip

\vfill

\begin{flushright}
    \footnotesize{
    Выполнил: студент 2-го курса магистратуры кафедры вычислительной физики, А. Р. Шабанов
    
    Руководитель: доцент кафедры вычислительной физики, к.ф.-м.н., А. С. Немнюгин
    
    Рецензент: разработчик интеллектуальных систем в ООО “Мэйл.ру”, А. А. Дзюба
    }
\end{flushright}

\end{center}

\vfill

\newlength{\ML}
\settowidth{\ML}{«\underline{\hspace{0.7cm}}» \underline{\hspace{2cm}}}


\begin{center}
    Санкт - Петербург \\
    2019
\end{center}

\end{titlepage}



\tableofcontents
\setcounter{page}{2}


\newpage
\indent
\indent
С каждым днем пользователи социальных сетей создают и потребляют все б\'{о}льшие
и б\'{о}льшие объемы информации, в том числе огромное количество фотографий.
Построение системы для быстрой и точной навигации в миллионах изображений
 --- не тривиальная задача. Одно из самых распространенных решений этой
 проблемы заключается в использовании тегов. Частный случай такого подхода ---
 использование хештегов в социальных сетях.


\indent
\textit{Хештег} --- это любое слово или фраза без пробелов, перед которой стоит 
символ \#, который называется \textit{диез} или \textit{решетка}, а в англоязычном 
варианте -- \textit{hash}, отсюда и название. Приведем несколько примеров:
\htag{masterwork}, \htag{spbu}, \#htag{lovesport}. Обычно в браузере или 
приложении хештег отображается как гипертекст, кликнув по которому можно 
получить список публикаций, снабженных таким же тегом.

\indent
Кроме простоты и удобства теги обладают еще одним полезным свойством 
-- они позволяют не думать об  иерархии структурируемой информации. 
Например, набор изображений можно разложить
по папкам, создав иерархию по датам, геолокациям или авторству. Причем в отдельных
случаях подобрать наиболее подходящую иерархию бывает затруднительно. Проблему
можно решить так: достаточно поставить несколько тегов для всех изображений,
а сами они могут храниться в плоской системе файлов.
Благодаря этому свойству тегирование используются  для рубрикации контента 
не только только онлайн, но и в оффлайн приложениях, например, 
просмоторщиках фотографий. 


\indent
Можно выделить две разновидности популярных хэштегов в социальных сетях.
Первые, используемые недолго и посвященные каким-то социальным явлениям
или событиям, например: \htag{elections2018}, \htag{metoo}. 
И вторые, широкораспространенные, но не связанные с новостной повесткой, 
например: \htag{sport}, \htag{cafe}; они и будут нас интересовать.
 Данная работа посвящена разработке интеллектуальной системы, 
подсказывающей пользователю релевантные хештеги к загружаемым фотографиям.
 Кроме того, с помощью такой системы можно решать и 
  \textquotedblleft обратную\textquotedblright  задачу -- определять,
уместно ли поставлены те или иные теги к заданным изображениям. 
Способность системы давать ответ на такой вопрос можно использовать для 
выявления злоупотреблений со стороны пользователей. Например, зачастую
 в рекламных целях продвигаемые публикации снабжают множеством популярных 
 тегов, не имеющих никакого отношения к публикуемой информации.
 
  
\addcontentsline{toc}{section}{\tocsecindent{Введение}}


\newpage
\section{Постановка задачи}

todo Рассказать про другие работы о тегах (например, на фликре).

\indent
\indent
Данная работа посвящена разработке интеллектуальной системы, 
подсказывающей пользователю релевантные и широкораспространенные
хештеги к загружаемым фотографиям.
Можно выделить две разновидности популярных хэштегов в социальных сетях.
Первые являются выражением чувств, эмоций и других абстрактных понятий,
например \htag{love}, \htag{friendship}, \htag{happy}. Вторые непосредственно
связаны с объектами, находящимися в кадре, например  
\htag{sport}, \htag{cafe}, \htag{nature}. 
Интуитивно понятно, что для второй разновидности
тегов предсказательная система, построенная на машинном зрении, будет
давать более точные результаты, чем для первой (но будет
иметь относительно узкую специализацию). Чтобы построить такую систему
нужен подходящий набор обучающих данных, но существующие датасеты, 
собранные из социальных сетей, не разделяют тренировочные примеры на те, 
где визуальный контент напрямую связан с целевыми классами, и те, где это не так.
В связи с чем возникла идея для данной задачи 
адаптировать датасет, изначально не размеченный
хэштегами, но имеющий аннотации, связанные с объектами,
находящимися в кадре. Кроме того, датасет должен быть относительно небольшим,
чтобы не требовать значительных 
вычислительных ресурсов; но при этом достаточно разнообразным, чтобы 
быть релевантным как можно большему количеству загружаемого контента.
В настоящей работе автор остановил свой выбор на датасете мест и локаций 
 \textit{SUN (Scene Understanding Dataset\cite{sundata} (SUN)}, 
 разметка которого после некоторой предобработки может быть отображена в популярные хэштеги.


\indent  
\textit{SUN} -- это набор изображений для каждого из которых выбрана одна
  из четырехсот локаций (сцен), вот несколько примеров:
  
  
\begin{itemize}
    \item \textit{baseball field}
    \item \textit{basketball court}
    \item \textit{ice shelf}
    \item \textit{forest}
    \item \textit{wind farm}
\end{itemize}


\indent
  Большинство названий локаций сами по себе не являются популярными 
  тегами из социальных сетей. Поэтому необходимо
  сопоставить их широко распространенным хэштегам (если это возможно):
  
  \begin{itemize}
      \item \textit{baseball field, basketball court}  $\rightarrow$ \htag{sport}
      \item \textit{ice shelf, forest} $\rightarrow$ \htag{nature}
      \item \textit{wind farm} $\rightarrow \times$ 
  \end{itemize}
  
  \indent
  Обученная на таких данных модель по окружению,
   обнаруженному на пользовательском фото,
  сможет подсказывать ему подходящий хэштег.
  
  
  
  \indent
  Ясно, что полученная модель будет корректно работать лишь для
  ограниченного (пусть и большого) домена фотографий. Следовательно, необходимо 
  обучить её распознавать не входящие в этот домен классы и
   не пытаться определить их категорию. Кроме того, в случае низкой уверенности
   в правильности предсказания так же лучше ничего не делать. 
    По мнению автора, гораздо предпочтительнее не 
   предложить пользователю подходящий хэштег, чем многократно предлагать 
   нерелевантные варианты.
   
   \indent
В качестве модели компьютерного зрения в работе 
используются глубокие сверточные
нейронные сети. Данный выбор обосновывается тем, что в последние
годы сверточные сети отлично зарекомендовали себя для решения задач, связанных
с обработкой изображений. Mожно привести в пример самый большой 
и известный
конкурс по классификации изображений \textit{ImageNet}\cite{imagenet},
который проводится ежегодно с 2010 года, при этом, начиная с 2012 года классические
решения ни разу не побеждали, уступив место нейронным сетям.


  \indent
Научная новизна работы определяется:
\begin{itemize}

    \item Адаптацией датасета \textit{SUN} для решения задачи о структурировании
    изображений в социальных сетях;
    
    \item Исследованием новых схем обучения неронных сетей, предложенных автором.
     (todo Обучить сегментатор на подсете, сделать псевдолейблинг, 
     посмотреть точность, 
    что-то вроде supervised atention)
    
\end{itemize}









\newpage
\section{Используемые инструменты}

\indent
\indent
\textbf{Программные средства}

\indent
В качестве языка программирования использовался
 \textit{Python 3.6.7} (сборка \textit{Anaconda}),
 в качестве среды разработки --- \textit{PyCharm Professional 2018.1},
 операционная система \textit{Ubuntu 16.04.4 LT}. Использованы
 следующие сторонние библиотеки для \textit{python}:


\begin{itemize}

    \item \textit{pytorch, torchvision} --- построение и обучение нейронных сетей
    \item \textit{tensorboardX} --- визуальное логирование процесса обучения 
    \item \textit{PIL, opencv, scikit-image, scipy} --- обработка изображений
    \item \textit{matplotlib} --- отрисовка графиков
    \item \textit{numpy} --- матричные вычисления
    \item \textit{pandas} --- работа с таблицами
    \item \textit{nltk} --- работа с текстом, в том числе с базой \textit{WordNet}
    \item \textit{scikit-learn} --- библиотека машинного обучения общего назначения
    \item \textit{pip} --- пакетный мененджер
    \item \textit{InstaLooter} --- утилита для автоматического скачивания
    изображений и видео из сети \textit{Instagram} по заданному хэштегу.
   
\end{itemize}


\bigskip

\indent
\indent
\textbf{Вычислительные мощности}

\indent
\indent
Обучение моделей производилось на удаленном сервере со следующей конфигурацией:
\begin{itemize}
    \item видеокарта \textit{GEFORCE GTX 1080 Ti} (\textit{11 ГБ} видеопамяти)
    \item процессор \textit{AMD Ryzen Threadripper 1920X 12-Core}
    \item оперативная память объемом  \textit{64 ГБ}
\end{itemize}



\newpage
\section{Сверточные нейронные сети}

\subsection{Введение в теорею сверточных нейросетей}

\subsection{Обзор используемых архитектур}


\section{Работа с данными}

\subsection{Датасет SUN}

\subsection{База WordNet}

\subsection{Адоптация датасета SUN}


\section{Численные эксперименты}

\subsection{Архитектура программы}

\subsection{Особенности реализации}

\subsection{Полученные результаты}

\subsection{Обсуждение результатов}



\section{Валидация результатов}



\section*{Выводы}
\addcontentsline{toc}{section}{\tocsecindent{Выводы}}


\newpage
\addcontentsline{toc}{section}{\tocsecindent{Список литературы}}
\begin{thebibliography}{99}

	\bigskip
	
	\bibitem{imagenet}
	Olga Russakovsky, Jia Deng, Hao Su, Jonathan Krause, Sanjeev Satheesh, 
	Sean Ma, Zhiheng Huang, Andrej Karpathy, Aditya Khosla, Michael Bernstein, 
	Alexander C. Berg and Li Fei-Fei. [2015] \textit{ImageNet Large Scale
	 Visual Recognition Challenge}. IJCV.
	 % http://ai.stanford.edu/~olga/bibtex/ILSVRC15.bib
	 
	 
	
\end{thebibliography}


\end{document} 
